\section{The nonintegral compact exponent}

\begin{proof}[Proof of Theorem~\ref{thm:main}]
Section~\ref{sec:geometry} constructs a smooth projective
Calabi--Yau threefold with torsion-free second homology and proves
primitivity of $\alpha=3s+4f$. The stable-pair argument of
Section~\ref{sec:pairs} gives \eqref{eq:primitive-count}.
Proposition~\ref{prop:congruence} now yields
\[
 e_{3\alpha,10}
 =c_{3\alpha,10}-\frac{4720}{3}(-14)
 =c_{3\alpha,10}+\frac{66080}{3}.
\]
The first term is an integer by Ionel--Parker's elementary
expansion, while $66080/3=22026+2/3$. This proves the theorem.
\end{proof}

The class $\alpha$ has divisor-degree vector
$(0,3,4,0,0,0)$ in \eqref{eq:degrees}; the nonintegral exponent
occurs at $(0,9,12,0,0,0)$ and local genus ten.
Corollary~\ref{cor:small-divisibility} shows that divisibility three
is the first possible divisibility for such a counterexample.
We do not claim that genus ten is the smallest possible local genus.

The rational factorization remains well-defined and unique.
Its arithmetic depends on the chosen local factors: the integer
elementary expansion and the integral geometric-generator algebra
are different structures. Formula~\eqref{eq:obstruction-intro}
isolates one condition that an integral scalar expansion would
impose on compact curve counts, and the surface calculation shows
that this condition fails.
